\begin{itemize}[leftmargin=*,itemsep=0.9\baselineskip,topsep=0.5\baselineskip]
  \item \textbf{Konstante zum Funktionswert addieren:}
    \par\smallskip
    \begin{minipage}[c]{0.40\linewidth}
      \raggedright
      \[
        g_C(x)=f(x)+C
      \]
      Bei $C=1$ wandert der Gipfel von $(0,1)$ nach $(0,2)$;
      alle Höhen steigen um $1$.
      \par\smallskip
      \emph{Folgerung:} $C$ verschiebt den Graphen vertikal um $C$.
    \end{minipage}\hfill
    \begin{minipage}[c]{0.55\linewidth}
      \centering
      \CachedTikzFigure{content/resources/chapters/real_functions/tikz/functions_gaussian_output_shift.tex}
    \end{minipage}

  \item \textbf{Konstante zum Argument addieren:}
    \par\smallskip
    \begin{minipage}[c]{0.40\linewidth}
      \raggedright
      \[
        g_C(x)=f(x+C)
      \]
      Bei $C=1$ liegt der Gipfel bei $(-1,1)$, bei $C=-1$ bei $(1,1)$.
      \par\smallskip
      \emph{Folgerung:} Alle $x$-Koordinaten ändern sich um $-C$.
    \end{minipage}\hfill
    \begin{minipage}[c]{0.55\linewidth}
      \centering
      \CachedTikzFigure{content/resources/chapters/real_functions/tikz/functions_gaussian_input_shift.tex}
    \end{minipage}

  \item \textbf{Funktionswert mit einer Konstanten multiplizieren:}
    \par\smallskip
    \begin{minipage}[c]{0.40\linewidth}
      \raggedright
      \[
        g_C(x)=C\,f(x)
      \]
      $C=2$ verdoppelt, $C=\tfrac12$ halbiert jede Höhe.
      \par\smallskip
      \emph{Folgerung:} Vertikaler Faktor $|C|$;
      bei $C<0$ zusätzlich Spiegelung an der $x$-Achse.
    \end{minipage}\hfill
    \begin{minipage}[c]{0.55\linewidth}
      \centering
      \CachedTikzFigure{content/resources/chapters/real_functions/tikz/functions_gaussian_output_scale.tex}
    \end{minipage}

  \item \textbf{Argument mit einer Konstanten multiplizieren:}
    \par\smallskip
    \begin{minipage}[c]{0.40\linewidth}
      \raggedright
      \[
        g_C(x)=f(Cx),\qquad C\ne0
      \]
      Bei $C=2$ wandert $(1,e^{-1})$ nach $(\tfrac12,e^{-1})$.
      \par\smallskip
      \emph{Folgerung:} Horizontaler Faktor $1/|C|$;
      bei $C<0$ zusätzlich Spiegelung an der $y$-Achse.
    \end{minipage}\hfill
    \begin{minipage}[c]{0.55\linewidth}
      \centering
      \CachedTikzFigure{content/resources/chapters/real_functions/tikz/functions_gaussian_input_scale.tex}
    \end{minipage}
\end{itemize}

\par\medskip
\begin{TheoremBox}{Verschieben, Strecken und Spiegeln}{real_functions_transformations}
  Die vier zuvor untersuchten \emph{Transformationen}
  \EnglishTerm{transformations} wirken für reelle Konstanten $C$ wie folgt:
  \begin{center}
    \renewcommand{\arraystretch}{1.5}
    \begin{tabular}{@{}lp{10.3cm}@{}}
      \toprule
      Transformation & Wirkung\\ \midrule
      $f(x)+C$ & vertikale Verschiebung um $C$\\
      $f(x+C)$ & horizontale Verschiebung um $-C$\\
      $C\,f(x)$ & vertikaler Faktor $|C|$;
        bei $C<0$ zusätzlich Spiegelung an der $x$-Achse\\
      $f(Cx)$, $C\ne0$ & horizontaler Faktor $1/|C|$;
        bei $C<0$ zusätzlich Spiegelung an der $y$-Achse\\
      \bottomrule
    \end{tabular}
  \end{center}
  Ein Faktor größer als $1$ streckt, ein Faktor zwischen $0$ und $1$ staucht
  in der jeweiligen Richtung. Bei $C=0$ ist $C\,f(x)$ die Nullfunktion.
\end{TheoremBox}
\begin{ExampleBox}{Eine Parabel schrittweise transformieren}{real_functions_transformed_parabola}
  Aus der bekannten Parabel $f_0(x)=x^2$ konstruieren wir
  $g(x)=-2(x-1)^2+3$ in drei Schritten:
  \begin{enumerate}[label=\arabic*.,leftmargin=*,itemsep=0.3\baselineskip,topsep=0.4\baselineskip]
    \item $f_1(x)=f_0(x-1)=(x-1)^2$: um $1$ nach rechts verschieben.
    \item $f_2(x)=-2f_1(x)=-2(x-1)^2$: an der $x$-Achse spiegeln
      und vertikal mit Faktor $2$ strecken.
    \item $g(x)=f_2(x)+3=-2(x-1)^2+3$: um $3$ nach oben verschieben.
  \end{enumerate}
  \par\medskip
  \IfClozeTF{Zeichnen Sie die drei Zwischen- und Ergebnisgraphen in die freien
    Koordinatensysteme.}{Blau zeigt den aktuellen Graphen, grau gestrichelt
    jeweils den Graphen vor dem letzten Schritt.}
  \par\smallskip
  {\centering
    \CachedTikzFigure{content/resources/chapters/real_functions/tikz/functions_transformed_parabola_steps.tex}\par}
  \par\medskip
\end{ExampleBox}

\par\medskip
Vor dem Ablesen einer horizontalen Verschiebung muss der innere Term
passend ausgeklammert werden:
$f(2x-4)=f(2(x-2))$ verschiebt den Graphen um $2$ nach rechts, nicht um $4$.